% Vietnamese translation for OI Public Library
% Source-File: IOI中国国家候选队论文/2005/李羽修/李羽修.doc
\documentclass[11pt]{article}
\usepackage{oipl-vn}

\newcommand{\Oh}{\mathcal{O}}
\newcommand{\floor}[1]{\left\lfloor #1 \right\rfloor}
\newcommand{\fracpart}[1]{\{#1\}}

\title{Thiết kế và tối ưu hàm băm}
\author{Li Yuxiu\\Trường Trung học Nankai, Thiên Tân}
\date{2005}

\begin{document}
\maketitle

\origtitle{Thiết kế và tối ưu hàm băm}
\sourcepath{IOI China candidate team papers / 2005 / Li Yuxiu.doc}

\renewcommand{\abstractname}{Tóm tắt}
\begin{abstract}
Băm là một cấu trúc dữ liệu thường dùng trong thi tin học. Một hàm băm tốt có
thể cải thiện đáng kể hiệu quả thời gian và bộ nhớ của toàn bộ chương trình.
Bài viết thảo luận cách thiết kế hàm băm cho nhiều loại mẫu khác nhau, tức các
giá trị khóa, đồng thời phân tích và tổng kết một số hàm băm thường dùng.
\end{abstract}

\noindent\textbf{Từ khóa:} băm; hàm băm; chuỗi; số nguyên; số thực; hoán vị và
tổ hợp.

\section{Tiêu chí đánh giá}

Với một hàm băm, tiêu chí quan trọng để đánh giá tốt hay xấu là tính ngẫu
nhiên: với một tập mẫu bất kỳ, xác suất rơi vào từng ô của bảng băm có đều
hay không. Xác suất càng đều thì dữ liệu phân bố trong bảng càng đều và hiệu
suất dùng bộ nhớ càng cao. Trong thi đấu, tính chất của mẫu thường không thể
biết trước, nên suy luận toán học bị hạn chế. Bài viết dùng thực nghiệm để
quan sát tính ngẫu nhiên ấy.

\section{Hàm băm cho số nguyên}

Ba phương pháp thường gặp là lấy dư trực tiếp, nhân rồi lấy phần nguyên, và
lấy giữa bình phương. Giả sử khóa là $k$, dung lượng bảng băm là $m$, và hàm
băm là $h(k)$.

\subsection{Lấy dư trực tiếp}

Ta chia khóa $k$ cho $m$ và dùng phần dư làm vị trí trong bảng:
\[
  h(k)=k\bmod m.
\]
Ví dụ, nếu dung lượng bảng là $701$ và khóa là $1234$, vị trí là
$1234\bmod 701=533$.

Ưu điểm của phương pháp này là dễ cài đặt, chạy nhanh, nên rất thường dùng.
Nhược điểm là nếu $m$ chọn không tốt và mẫu lại có cấu trúc đặc biệt, dữ liệu
dễ dồn vào một số ô. Theo kinh nghiệm, $m$ thường nên là một số nguyên tố
không quá gần một lũy thừa đặc biệt; nếu miền giá trị của khóa nhỏ, có thể
chọn $m$ trong khoảng $1.1$ tới $1.6$ lần miền giá trị. Khi thi đấu, có thể
viết bộ sinh số nguyên tố hoặc chọn một số ``trông giống nguyên tố''.

Tác giả thử chèn $4000$ số vào bảng dung lượng $701$:
\[
\begin{array}{lrr}
\text{Dữ liệu} & \text{ngẫu nhiên} & \text{liên tiếp}\\
\hline
\text{ô nhỏ nhất} & 0 & 5\\
\text{ô lớn nhất} & 15 & 6\\
\text{kỳ vọng} & 5.70613 & 5.70613\\
\text{độ lệch chuẩn} & 2.4165 & 0.455531
\end{array}
\]
Với dữ liệu ngẫu nhiên, ô lớn nhất gần gấp ba kỳ vọng. Trên máy Pentium III
866MHz, 128MB RAM của tác giả, một lần tính mất khoảng $39$ ns, tức khoảng
$35$ chu kỳ.

\subsection{Nhân rồi lấy phần nguyên}

Chọn một số thực $A\in(0,1)$, tốt nhất là vô tỉ. Tính $kA$, lấy phần thập
phân, nhân với $m$, rồi lấy phần nguyên:
\[
  h(k)=\floor{m\fracpart{kA}}.
\]
Ở đây $\fracpart{x}=x-\floor{x}$. Ví dụ có thể chọn $A=(\sqrt5-1)/2$, một lựa
chọn thường cho hiệu quả tốt.

Với $4000$ số và bảng dung lượng $701$, tác giả ghi nhận:
\[
\begin{array}{lrr}
\text{Dữ liệu} & \text{ngẫu nhiên} & \text{liên tiếp}\\
\hline
\text{ô nhỏ nhất} & 0 & 4\\
\text{ô lớn nhất} & 15 & 7\\
\text{kỳ vọng} & 5.70613 & 5.70613\\
\text{độ lệch chuẩn} & 2.5069 & 0.619999
\end{array}
\]
Từ công thức có thể thấy phương pháp này ít phụ thuộc vào $m$, nên khi $m$
không phù hợp cho lấy dư trực tiếp, nó có thể biểu hiện tốt. Tuy nhiên trong
thử nghiệm trên, kết quả không quá lý tưởng; hơn nữa do dùng nhiều phép toán
dấu phẩy động, tốc độ chậm. Sau nhiều lần tối ưu, tác giả vẫn cần khoảng
$892$ ns cho một lần tính, tức khoảng $810$ chu kỳ, chậm hơn lấy dư trực tiếp
khoảng $23$ lần.

\subsection{Lấy giữa bình phương}

Ta bình phương khóa $k$, rồi lấy một số bit hoặc chữ số ở giữa làm giá trị
băm. Vì mỗi chữ số của $k$ đều ảnh hưởng tới phần giữa của $k^2$, phương pháp
này thường có hiệu quả khá tốt. Tuy nhiên với $k$ nhỏ, hiệu quả không lý
tưởng; cài đặt cũng rườm rà hơn. Để tận dụng không gian bảng băm, dung lượng
$m$ nên là lũy thừa của $2$.

Với $4000$ số và bảng dung lượng $512$, tác giả ghi nhận:
\[
\begin{array}{lrr}
\text{Dữ liệu} & \text{ngẫu nhiên} & \text{liên tiếp}\\
\hline
\text{ô nhỏ nhất} & 0 & 1\\
\text{ô lớn nhất} & 17 & 17\\
\text{kỳ vọng} & 7.8125 & 7.8125\\
\text{độ lệch chuẩn} & 2.95804 & 2.64501
\end{array}
\]
Kết quả kém hơn tưởng tượng, nhất là với dữ liệu liên tiếp. Nhưng do chỉ dùng
phép nhân và phép bit, nó rất nhanh: khoảng $23$ ns, tức $19$ chu kỳ, còn
nhanh hơn lấy dư trực tiếp.

\subsection{So sánh}

\[
\begin{array}{llll}
\text{Phương pháp} & \text{Cài đặt} & \text{Hiệu quả} & \text{Tốc độ}\\
\hline
\text{lấy dư trực tiếp} & \text{dễ} & \text{tốt} & \text{khá nhanh}\\
\text{nhân lấy phần nguyên} & \text{khá dễ} & \text{khá tốt} & \text{chậm}\\
\text{lấy giữa bình phương} & \text{trung bình} & \text{khá tốt} & \text{nhanh}
\end{array}
\]
Từ bảng này có thể thấy lấy dư trực tiếp có tỷ lệ hiệu quả trên độ phức tạp
tốt nhất, nên cũng là phương pháp thường dùng nhất trong thi đấu.

Với số thực, có thể trực tiếp dùng phương pháp nhân rồi lấy phần nguyên. Với
những loại dữ liệu khác, ta thường chuyển chúng thành số nguyên trước rồi mới
đưa vào bảng băm.

\section{Hàm băm cho chuỗi}

Bản thân một chuỗi có thể xem là một số nguyên lớn trong cơ số $256$; với
chuỗi ANSI có thể xem là cơ số $128$. Vì vậy có thể dùng lấy dư trực tiếp và
tính giá trị trong thời gian tuyến tính theo độ dài chuỗi. Để đảm bảo hiệu
quả, $m$ vẫn không nên chọn quá gần các giá trị đặc biệt. Đặc biệt, nếu xem
chuỗi là số trong cơ số $b$, chọn $m=b-1$ sẽ làm mọi hoán vị của cùng một
chuỗi có cùng giá trị băm, vì giá trị theo modulo $b-1$ chỉ phụ thuộc vào tổng
chữ số.

Các hàm băm chuỗi thường dùng còn có ELFHash, APHash, v.v. Chúng dùng phép
toán bit để làm mỗi ký tự ảnh hưởng tới giá trị cuối. Ngoài ra còn có các hàm
băm mật mã như MD5 và SHA1, vốn rất khó tìm va chạm trong bối cảnh thông
thường.

Tác giả lấy ngẫu nhiên $1000$ từ khác nhau và $1000$ câu khác nhau từ một tiểu
thuyết của Mark Twain, lần lượt làm dữ liệu chuỗi ngắn và chuỗi dài. Sau khi
dùng các hàm băm biến chúng thành số nguyên, lại lấy dư trực tiếp vào bảng
dung lượng $1237$; nếu xung đột thì chuỗi mới ghi đè chuỗi cũ. Số ô còn được
phủ sau cùng phản ánh thô hiệu quả của hàm băm:
\[
\begin{array}{lrrr}
\text{Hàm} & \text{ngắn} & \text{dài} & \text{trung bình}\\
\hline
\text{lấy dư trực tiếp} & 667 & 676 & 671.5\\
\text{P. J. Weinberger} & 683 & 676 & 679.5\\
\text{ELF} & 683 & 676 & 679.5\\
\text{SDBM} & 694 & 680 & 687.0\\
\text{BKDR} & 665 & 710 & 687.5\\
\text{DJB} & 694 & 683 & 688.5\\
\text{AP} & 684 & 698 & 691.0\\
\text{RS} & 691 & 693 & 692.0\\
\text{JS} & 684 & 708 & 696.0
\end{array}
\]
Đưa $1000$ số ngẫu nhiên vào bảng dung lượng $1237$ bằng lấy dư trực tiếp cũng
chỉ phủ được $694$ ô, nên vài phương pháp phía sau đã gần tới giới hạn ngẫu
nhiên. Tác giả thường chọn JS Hash hoặc SDBM Hash. Hai đoạn mã được trích lại
nguyên văn vì chúng là mã C chuẩn, không phải nội dung cần dịch:

\begin{verbatim}
unsigned int JSHash(char *str)
{
    unsigned int hash = 1315423911;
    while (*str)
        hash ^= ((hash << 5) + (*str++) + (hash >> 2));
    return (hash & 0x7FFFFFFF);
}

unsigned int SDBMHash(char *str)
{
    unsigned int hash = 0;
    while (*str)
        hash = (*str++) + (hash << 6) + (hash << 16) - hash;
    return (hash & 0x7FFFFFFF);
}
\end{verbatim}

JSHash có phép tính phức tạp hơn; nếu yêu cầu hiệu quả không quá khắt khe,
SDBMHash là một lựa chọn tốt.

\section{Hàm băm cho hoán vị}

Chính xác hơn, phần này không chỉ là ``hashing'', mà là một quá trình đánh số
(\textit{numerize}): xây dựng song ánh giữa các hoán vị và các số tự nhiên.
Khi đó ta có thể định địa chỉ trực tiếp, dùng làm hàm băm, hoặc dùng trong
quy hoạch động nén trạng thái.

\subsection{Kiến thức nền}

Ta quen với biểu diễn cơ số $b$:
\[
  x=a_0+a_1b+a_2b^2+\cdots.
\]
Tương tự, dùng hệ cơ số biến đổi dựa trên giai thừa. Từ đẳng thức
\[
  n!=(n-1)(n-1)!+(n-1)!,
\]
có thể suy ra mọi số nguyên từ $0$ tới $n!-1$ đều biểu diễn duy nhất được dưới
dạng
\[
  x=a_1\cdot 1!+a_2\cdot 2!+\cdots+a_{n-1}\cdot (n-1)!,
\]
với
\[
  0\le a_i\le i.
\]
Có thể hiểu dãy $a_{n-1},\ldots,a_1$ là một ``số cơ số biến đổi'': chữ số thứ
nhất có cơ số $2$, chữ số thứ hai có cơ số $3$, v.v. Vì có đúng $n!$ dãy như
vậy, ta tự nhiên liên hệ chúng với $n!$ hoán vị của $n$ phần tử.

\subsection{Song ánh giữa hoán vị đầy đủ và số tự nhiên}

Giả sử các phần tử là $1,2,\ldots,n$. Với một hoán vị $p$, đặt $a_i$ là số
phần tử nhỏ hơn $i$ nằm bên phải $i$ trong hoán vị. Dãy các $a_i$ chính là mã
Lehmer theo giá trị phần tử.

Ví dụ, hoán vị $4213$ của $1,2,3,4$ có: bên phải $4$ có $3$ số nhỏ hơn, bên
phải $3$ có $0$ số nhỏ hơn, bên phải $2$ có $1$ số nhỏ hơn. Vì vậy mã là
$301$, ứng với
\[
  3\cdot 3!+0\cdot 2!+1\cdot 1!=19.
\]
Ngược lại, từ $19$ có thể khôi phục mã $301$, rồi khôi phục hoán vị $4213$.

\subsection{Hoán vị chọn một phần}

Ta cũng có thể đánh số các chỉnh hợp chọn $n$ phần tử từ $m$ phần tử. Theo
nguyên lý nhân,
\[
  P(m,n)=m(m-1)(m-2)\cdots(m-n+1).
\]
Ở vị trí thứ nhất có $m$ cách chọn, vị trí thứ hai có $m-1$ cách, v.v. Vì vậy
ta định nghĩa một hệ cơ số biến đổi $m$--$n$: chữ số thứ nhất có cơ số $m$,
chữ số thứ hai có cơ số $m-1$, ..., chữ số thứ $n$ có cơ số $m-n+1$.

Mọi số $x$ trong khoảng $[0,P(m,n)-1]$ có thể biểu diễn duy nhất bằng dãy chữ
số tương ứng. Để ánh xạ một chỉnh hợp sang số, duy trì danh sách
\[
  L=[1,2,\ldots,m]
\]
đánh chỉ số từ $0$. Với chỉnh hợp $p_1,p_2,\ldots,p_n$, lần lượt tìm chỉ số
của $p_i$ trong $L$, ghi lại chỉ số đó rồi xóa $p_i$ khỏi $L$.

Ví dụ, chọn $2,4,3$ từ $1,2,3,4,5$. Ban đầu $L=[1,2,3,4,5]$; lấy $2$ có chỉ
số $1$, rồi $L=[1,3,4,5]$; lấy $4$ có chỉ số $2$, rồi $L=[1,3,5]$; lấy $3$ có
chỉ số $1$. Mã là $121$ trong hệ cơ số $5,4,3$, ứng với
\[
  1\cdot 4\cdot 3+2\cdot 3+1=19.
\]
Ngược lại, từ số $19$ có thể khôi phục mã $121$ và chỉnh hợp $243$.

\section{Kết luận}

Bài viết đã giới thiệu và phân tích một số hàm băm thường dùng, đồng thời mở
rộng sang quá trình xây dựng song ánh giữa cấu trúc tổ hợp và số tự nhiên.
Băm là một cấu trúc dữ liệu hiệu quả, thể hiện rõ tư tưởng dùng không gian đổi
lấy thời gian. Khi giới hạn bộ nhớ trong các kỳ thi ngày càng thoáng hơn, biết
tận dụng bộ nhớ để đổi lấy thời gian là một kỹ năng quan trọng; băm có thể
giúp rất nhiều trong việc đó. Cần căn cứ vào đặc điểm của đề để chọn cấu trúc
dữ liệu và hàm băm phù hợp.

Tác giả cũng nêu rằng với tổ hợp chưa tìm được một cách song ánh tốt như với
hoán vị, và mong được thảo luận thêm.

\section*{Tài liệu tham khảo}

\begin{itemize}
  \item Thomas H. Cormen, Charles E. Leiserson, Ronald L. Rivest, Clifford
  Stein, \textit{Introduction to Algorithms}, Second Edition, MIT Press, 2001.
  \item Liu Rujia, Huang Liang, \textit{Nghệ thuật thuật toán và thi tin học},
  Nhà xuất bản Đại học Thanh Hoa, 2004.
  \item Lu Kaicheng, Lu Huaming, \textit{Toán tổ hợp}, ấn bản thứ ba, Nhà xuất
  bản Đại học Thanh Hoa, 2002.
\end{itemize}

\appendix
\section{Ghi chú về phụ lục mã nguồn}

Bản gốc kèm mã C của nhiều hàm băm chuỗi thường dùng: RS, JS, P. J.
Weinberger, ELF, BKDR, SDBM, DJB và AP. Vì đây là các đoạn mã mẫu và yêu cầu
hiện tại không cần dịch code templates, phụ lục này chỉ ghi nhận danh sách
hàm. Hai hàm JS và SDBM đã được trích trong thân bài vì chúng là ví dụ trực
tiếp được tác giả khuyến nghị.

\end{document}
